\section{Discussion}
\label{sec:discussion}

\subsection{Interpretation of the controlled demonstration}

The implemented state machine enforces four useful properties in the frozen
examples. It retains negative method decisions, connects a successor to its
parent result, blocks post-outcome changes to the current contract, and
requires report claims to resolve to recorded evidence. The component
removals behave consistently with the co-designed fault rules, and exact
reruns recover the same cell hashes. These observations support an engineering
claim about enforcement in this package.

They do not establish an effect on a broader task population. The ten tasks
were constructed alongside the controller and evaluator, the repair policy is
deterministic, and three repetition identifiers duplicate each task outcome.
The additive analysis corrects the statistical unit but cannot create new
independent evidence. Its wide [\TaskLevelCILower{}, \TaskLevelCIUpper{}]
interval and two-sided sign-test \(p=\SignTestTwoSided{}\) are compatible with
a useful local mechanism test and insufficient for a comparative performance
claim.

The two negative equation-discovery rounds strengthen a different part of the
case. They show that the workflow can keep favorable development results and
unfavorable unseen decisions in one lineage without editing the acceptance
rule. They also prevent the Route B artifact from being read as evidence for
a new equation-discovery method.

\subsection{Position relative to current systems}

The current peer-reviewed field already contains end-to-end machine-learning
automation, multi-agent hypothesis development, laboratory-linked iterative
analysis, and large-scale empirical software search
~\cite{lu2026aiscientist,gottweis2026coscientist,ghareeb2026robin,
aygun2026era}. The present work is smaller in scope. Its possible value lies
in a testable state contract: when failure becomes a stored parent, when a new
round counts, which data can cross the outcome boundary, and which claim must
be blocked.

Independent benchmark evidence also changes the appropriate comparison.
AstaBench and REPRO-Bench show that broad scientific assistance and realistic
reproducibility assessment remain difficult even when evaluation is external
to the tested agent~\cite{bragg2026astabench,hu2025reprobench}. The current
matrix does not compare AutoResearch with those agents under common budgets.
It should therefore be described as a mechanism and integration test, not a
ranking result.

\subsection{Why the present object is not submission-ready}

Four missing layers explain the gap between a reproducible local object and a
publication-level systems result. First, the hash graph is project-specific;
an interoperable RO-Crate and PROV-O overlay with rights and software identity
is still required. Second, the field map and benchmark choice need independent
human coding rather than author selection. Third, a new confirmation must use
independently authored tasks from several substantive families, stochastic
controller trajectories where applicable, external systems and simple
baselines under common budgets, null controls, a prospective task-level power
analysis, and an independent scorer. Fourth, scientific validity, novelty,
ethics, authorship, licensing, disclosure, venue fit, and public distribution
need independent human review and adjudication.

These layers are non-compensating. More artifact hashes cannot replace new
tasks; a new benchmark cannot replace rights review; and an author rerun
cannot replace independent validation. A future negative confirmation can be
scientifically useful if its diagnostic question and power are frozen before
outcome access. Otherwise, the correct product is a research record rather
than a forced paper claim.

\subsection{Implications for research-system design}

Evidence lineage should be represented as state rather than reconstructed
from chat history. A parent hash and declared mechanism delta make retries and
scientific successors distinguishable. Outcome access should be one-way for a
given contract. Claims and successful execution should remain separate
metrics because a completed workflow can still produce an unsupported report.
Finally, open-science metadata should describe both positive and negative
activities while keeping permissions and human responsibility explicit.
